\begin{PSbox}[unbreakable]{obj}{Learning Objectives}

After completing this module, students will be able to:

\begin{enumerate}
\def\labelenumi{\arabic{enumi}.}
\tightlist
\item
  Formulate null and alternative hypotheses for engineering problems
\item
  Identify Type I and Type II errors and their engineering consequences
\item
  Compute test statistics and p-values
\item
  Perform Z-tests and t-tests for means
\item
  Correctly interpret p-values (and avoid common misconceptions)
\item
  Choose between one-sided and two-sided tests
\end{enumerate}

\end{PSbox}

\PSsection{15.1}{Introduction: Making Decisions from Data}

\textbf{Engineering question:} ``Is this sensor biased?'' ``Did the process change?'' ``Does the new algorithm perform better?''

We need a \textbf{systematic framework} to answer such questions from data, while controlling the probability of making errors.

\PSsection{15.2}{Framework}

\PSsubsection{Null Hypothesis H\textsubscript{0}}

The \textbf{default assumption} — typically ``nothing changed'' or ``no effect'':

\begin{itemize}
\tightlist
\item
  H\textsubscript{0}: \ensuremath{\mu} = \ensuremath{\mu}\textsubscript{0} (sensor is not biased)
\item
  H\textsubscript{0}: \ensuremath{\mu}\textsubscript{1} = \ensuremath{\mu}\textsubscript{2} (two methods perform equally)
\item
  H\textsubscript{0}: \ensuremath{\sigma}\textsuperscript{2} = \ensuremath{\sigma}\textsubscript{0}\textsuperscript{2} (variance hasn’t changed)
\end{itemize}

\PSsubsection{Alternative Hypothesis H\textsubscript{1}}

What we want to \textbf{detect} or prove:

\begin{itemize}
\tightlist
\item
  H\textsubscript{1}: \ensuremath{\mu} \ensuremath{\neq} \ensuremath{\mu}\textsubscript{0} (sensor IS biased — two-sided)
\item
  H\textsubscript{1}: \ensuremath{\mu} \textgreater{} \ensuremath{\mu}\textsubscript{0} (performance improved — one-sided)
\item
  H\textsubscript{1}: \ensuremath{\sigma}\textsuperscript{2} \textgreater{} \ensuremath{\sigma}\textsubscript{0}\textsuperscript{2} (variance increased)
\end{itemize}

\PSsubsection{Decision}

Based on data, either:

\begin{itemize}
\tightlist
\item
  \textbf{Reject H\textsubscript{0}} \ensuremath{\to} evidence supports H\textsubscript{1}
\item
  \textbf{Fail to reject H\textsubscript{0}} \ensuremath{\to} insufficient evidence against H\textsubscript{0}
\end{itemize}

\begin{PSquote}

\PSwarn{} ``Fail to reject H\textsubscript{0}'' \ensuremath{\neq} ``H\textsubscript{0} is true''. Absence of evidence \ensuremath{\neq} evidence of absence.

\end{PSquote}

\PSsection{15.3}{Types of Errors}

{\def\LTcaptype{none} % do not increment counter
\begin{longtable}[c]{@{}ccc@{}}
\toprule\noalign{}
\rowcolor{PSnavy}& \PSth{H\textsubscript{0} True (reality)} & \PSth{H\textsubscript{1} True (reality)} \\
\midrule\noalign{}
\endhead
\bottomrule\noalign{}
\endlastfoot
\textbf{Reject H\textsubscript{0}} (decision) & Type I Error (\ensuremath{\alpha}) & Correct! (Power = 1-\ensuremath{\beta}) \\
\textbf{Fail to reject H\textsubscript{0}} (decision) & Correct! & Type II Error (\ensuremath{\beta}) \\
\end{longtable}
}

\PSsubsection{Type I Error (False Alarm) — probability \ensuremath{\alpha}}

Rejecting H\textsubscript{0} when it’s actually true. ``Crying wolf.''

Engineering analogy: Declaring a process has changed when it hasn’t \ensuremath{\to} unnecessary recalibration.

\PSsubsection{Type II Error (Missed Detection) — probability \ensuremath{\beta}}

Failing to reject H\textsubscript{0} when H\textsubscript{1} is actually true. ``Missing the signal.''

Engineering analogy: Failing to detect a sensor bias \ensuremath{\to} biased measurements continue.

\PSsubsection{Engineering Analogy: Radar Detection}

{\def\LTcaptype{none} % do not increment counter
\begin{longtable}[c]{@{}cc@{}}
\toprule\noalign{}
\rowcolor{PSnavy}\PSth{Statistical Term} & \PSth{Radar Equivalent} \\
\midrule\noalign{}
\endhead
\bottomrule\noalign{}
\endlastfoot
H\textsubscript{0}: no target & No aircraft present \\
H\textsubscript{1}: target present & Aircraft present \\
Type I error (\ensuremath{\alpha}) & False alarm \\
Type II error (\ensuremath{\beta}) & Missed detection \\
Power (1-\ensuremath{\beta}) & Probability of detection \\
\end{longtable}
}

\PSsubsection{The Tradeoff}

Reducing \ensuremath{\alpha} (fewer false alarms) \ensuremath{\to} increases \ensuremath{\beta} (more missed detections), and vice versa.

Fix \ensuremath{\alpha} (significance level), then maximize power = 1 - \ensuremath{\beta}.

\PSsection{15.4}{Significance Level \ensuremath{\alpha}}

The \textbf{maximum acceptable Type I error rate}, chosen BEFORE looking at data.

Common choices: \ensuremath{\alpha} = 0.05, 0.01, 0.10

\textbf{Choosing \ensuremath{\alpha}:} Consider the cost of a false alarm:

\begin{itemize}
\tightlist
\item
  \ensuremath{\alpha} = 0.01: Very conservative (e.g., medical device safety)
\item
  \ensuremath{\alpha} = 0.05: Standard in most engineering/science
\item
  \ensuremath{\alpha} = 0.10: More tolerant (exploratory analysis)
\end{itemize}

\PSsection{15.5}{Power of a Test}

\PSdisplay{\text{Power} = 1 - \beta = P(\text{reject } H_0 | H_1 \text{ is true})}

Power depends on:

\begin{enumerate}
\def\labelenumi{\arabic{enumi}.}
\tightlist
\item
  \textbf{Effect size:} Larger true difference \ensuremath{\to} easier to detect \ensuremath{\to} more power
\item
  \textbf{Sample size n:} More data \ensuremath{\to} more power
\item
  \textbf{Significance level \ensuremath{\alpha}:} Larger \ensuremath{\alpha} \ensuremath{\to} more power (but more false alarms)
\item
  \textbf{Variability \ensuremath{\sigma}:} Less noise \ensuremath{\to} more power
\end{enumerate}

\PSsection{15.6}{The p-value}

\begin{PSbox}[unbreakable]{def}{Definition}

The \textbf{p-value} is the probability of obtaining a test statistic \textbf{as extreme or more extreme} than the observed value, \textbf{assuming H\textsubscript{0} is true}.

\PSdisplay{p\text{-value} = P(\text{data this extreme or more} | H_0 \text{ true})}

\end{PSbox}

\PSsubsection{Decision Rule}

\begin{itemize}
\tightlist
\item
  If p-value \textless{} \ensuremath{\alpha} \ensuremath{\to} Reject H\textsubscript{0}
\item
  If p-value \ensuremath{\ge} \ensuremath{\alpha} \ensuremath{\to} Fail to reject H\textsubscript{0}
\end{itemize}

\PSsubsection{\PSyes{} CORRECT Interpretation}

\begin{PSquote}

``The p-value is the probability of seeing data this extreme (or more) if H\textsubscript{0} were true. A small p-value means the data is unlikely under H\textsubscript{0}.''

\end{PSquote}

\PSsubsection{\PSno{} WRONG Interpretation}

\begin{PSquote}

{``The p-value is the probability that H\textsubscript{0} is true.''}

\end{PSquote}

\textbf{Why wrong:} H\textsubscript{0} is either true or false (fixed). The p-value says nothing about P(H\textsubscript{0} true). It measures how surprising the data is under H\textsubscript{0}.

\PSsubsection{Calibration}

{\def\LTcaptype{none} % do not increment counter
\begin{longtable}[c]{@{}cc@{}}
\toprule\noalign{}
\rowcolor{PSnavy}\PSth{p-value} & \PSth{Evidence against H\textsubscript{0}} \\
\midrule\noalign{}
\endhead
\bottomrule\noalign{}
\endlastfoot
\textgreater{} 0.10 & Weak or none \\
0.05 – 0.10 & Marginal \\
0.01 – 0.05 & Moderate \\
0.001 – 0.01 & Strong \\
\textless{} 0.001 & Very strong \\
\end{longtable}
}

\PSsection{15.7}{One-Sided vs. Two-Sided Tests}

\PSsubsection{Two-Sided (H\textsubscript{1}: \ensuremath{\mu} \ensuremath{\neq} \ensuremath{\mu}\textsubscript{0})}

Used when deviation in \textbf{either} direction is important.

\begin{itemize}
\tightlist
\item
  Reject if \textbar{}T\textbar{} \textgreater{} t\textsubscript{\ensuremath{\alpha}/2, n-1}
\item
  p-value = 2\ensuremath{\cdot}P(T \textgreater{} \textbar{}t\textsubscript{obs}\textbar{})
\end{itemize}

\PSsubsection{One-Sided (H\textsubscript{1}: \ensuremath{\mu} \textgreater{} \ensuremath{\mu}\textsubscript{0} or H\textsubscript{1}: \ensuremath{\mu} \textless{} \ensuremath{\mu}\textsubscript{0})}

Used when only one direction of deviation matters.

\begin{itemize}
\tightlist
\item
  H\textsubscript{1}: \ensuremath{\mu} \textgreater{} \ensuremath{\mu}\textsubscript{0}: Reject if T \textgreater{} t\textsubscript{\ensuremath{\alpha}, n-1}
\item
  H\textsubscript{1}: \ensuremath{\mu} \textless{} \ensuremath{\mu}\textsubscript{0}: Reject if T \textless{} -t\textsubscript{\ensuremath{\alpha}, n-1}
\end{itemize}

\PSsubsection{When to Use Which}

{\def\LTcaptype{none} % do not increment counter
\begin{longtable}[c]{@{}cc@{}}
\toprule\noalign{}
\rowcolor{PSnavy}\PSth{Scenario} & \PSth{Test Type} \\
\midrule\noalign{}
\endhead
\bottomrule\noalign{}
\endlastfoot
Is the sensor biased (either direction)? & Two-sided \\
Does the new algorithm IMPROVE BER? & One-sided (H\textsubscript{1}: BER \textless{} BER\textsubscript{old}) \\
Has noise power INCREASED? & One-sided (H\textsubscript{1}: \ensuremath{\sigma}\textsuperscript{2} \textgreater{} \ensuremath{\sigma}\textsubscript{0}\textsuperscript{2}) \\
Did the process mean CHANGE? & Two-sided \\
\end{longtable}
}

\PSsection{15.8}{Tests for Means}

\PSsubsection{One-Sample Z-Test (\ensuremath{\sigma} known)}

Test: H\textsubscript{0}: \ensuremath{\mu} = \ensuremath{\mu}\textsubscript{0}

Test statistic: Z = (\ensuremath{\bar{X}} - \ensuremath{\mu}\textsubscript{0})/(\ensuremath{\sigma}/\ensuremath{\surd}n)

Under H\textsubscript{0}: Z \textasciitilde{} N(0,1)

\PSsubsection{One-Sample t-Test (\ensuremath{\sigma} unknown)}

Test: H\textsubscript{0}: \ensuremath{\mu} = \ensuremath{\mu}\textsubscript{0}

Test statistic: T = (\ensuremath{\bar{X}} - \ensuremath{\mu}\textsubscript{0})/(S/\ensuremath{\surd}n)

Under H\textsubscript{0}: T \textasciitilde{} t(n-1)

\PSsubsection{Two-Sample t-Test (comparing two means)}

Test: H\textsubscript{0}: \ensuremath{\mu}\textsubscript{1} = \ensuremath{\mu}\textsubscript{2}

Test statistic: T = (\ensuremath{\bar{X}}\textsubscript{1} - \ensuremath{\bar{X}}\textsubscript{2}) / \ensuremath{\surd}(S\textsubscript{1}\textsuperscript{2}/n\textsubscript{1} + S\textsubscript{2}\textsuperscript{2}/n\textsubscript{2})

Approximate df by Welch-Satterthwaite formula.

\PSsection{15.9}{Test for Variance (Chi-Square Test)}

\PSsubsection{One-Sample Test}

H\textsubscript{0}: \ensuremath{\sigma}\textsuperscript{2} = \ensuremath{\sigma}\textsubscript{0}\textsuperscript{2}

Test statistic: \ensuremath{\chi}\textsuperscript{2} = (n-1)S\textsuperscript{2}/\ensuremath{\sigma}\textsubscript{0}\textsuperscript{2}

Under H\textsubscript{0}: \ensuremath{\chi}\textsuperscript{2} \textasciitilde{} \ensuremath{\chi}\textsuperscript{2}(n-1)

\PSsubsection{Decision}

\begin{itemize}
\tightlist
\item
  H\textsubscript{1}: \ensuremath{\sigma}\textsuperscript{2} \textgreater{} \ensuremath{\sigma}\textsubscript{0}\textsuperscript{2}: Reject if \ensuremath{\chi}\textsuperscript{2} \textgreater{} \ensuremath{\chi}\textsuperscript{2}\textsubscript{\ensuremath{\alpha}, n-1}
\item
  H\textsubscript{1}: \ensuremath{\sigma}\textsuperscript{2} \ensuremath{\neq} \ensuremath{\sigma}\textsubscript{0}\textsuperscript{2}: Reject if \ensuremath{\chi}\textsuperscript{2} \textless{} \ensuremath{\chi}\textsuperscript{2}\textsubscript{1-\ensuremath{\alpha}/2, n-1} or \ensuremath{\chi}\textsuperscript{2} \textgreater{} \ensuremath{\chi}\textsuperscript{2}\textsubscript{\ensuremath{\alpha}/2, n-1}
\end{itemize}

\PSsection{15.10}{Step-by-Step Hypothesis Testing Procedure}

\begin{enumerate}
\def\labelenumi{\arabic{enumi}.}
\tightlist
\item
  \textbf{State hypotheses:} H\textsubscript{0} and H\textsubscript{1}
\item
  \textbf{Choose significance level:} \ensuremath{\alpha} (e.g., 0.05)
\item
  \textbf{Select test statistic:} Z, t, or \ensuremath{\chi}\textsuperscript{2} depending on scenario
\item
  \textbf{Determine critical value(s)} or compute p-value
\item
  \textbf{Compute test statistic} from data
\item
  \textbf{Make decision:} Reject H\textsubscript{0} if test statistic in critical region (or p \textless{} \ensuremath{\alpha})
\item
  \textbf{State conclusion} in engineering terms
\end{enumerate}

\PSsection{15.11}{Engineering Examples}

\begin{PSbox}[unbreakable]{ex}{Example 1: Is a Sensor Biased?}

A sensor should read 0V with no input. 20 measurements give \ensuremath{\bar{X}} = 0.012V, S = 0.03V.

H\textsubscript{0}: \ensuremath{\mu} = 0 (no bias), H\textsubscript{1}: \ensuremath{\mu} \ensuremath{\neq} 0 (biased), \ensuremath{\alpha} = 0.05

T = (0.012 - 0)/(0.03/\ensuremath{\surd}20) = 0.012/0.00671 = 1.789

t\textsubscript{0.025, 19} = 2.093. Since \textbar{}T\textbar{} = 1.789 \textless{} 2.093 \ensuremath{\to} \textbf{Fail to reject H\textsubscript{0}.}

p-value = 2\ensuremath{\cdot}P(t(19) \textgreater{} 1.789) \ensuremath{\approx} 0.090. Not significant at 5\% level.

Conclusion: Insufficient evidence to conclude the sensor is biased.

\end{PSbox}

\begin{PSbox}[unbreakable]{ex}{Example 2: Did Manufacturing Process Change?}

Before: \ensuremath{\mu}\textsubscript{0} = 100\ensuremath{\Omega} (established). After modification: n = 30, \ensuremath{\bar{X}} = 101.2\ensuremath{\Omega}, S = 3.5\ensuremath{\Omega}.

H\textsubscript{0}: \ensuremath{\mu} = 100, H\textsubscript{1}: \ensuremath{\mu} \ensuremath{\neq} 100, \ensuremath{\alpha} = 0.05

T = (101.2 - 100)/(3.5/\ensuremath{\surd}30) = 1.2/0.639 = 1.878

t\textsubscript{0.025, 29} = 2.045. \textbar{}T\textbar{} = 1.878 \textless{} 2.045 \ensuremath{\to} \textbf{Fail to reject H\textsubscript{0}.}

\end{PSbox}

\begin{PSbox}[unbreakable]{ex}{Example 3: Does New Algorithm Improve BER?}

Old algorithm: BER\textsubscript{0} = 10\textsuperscript{-3}. New algorithm tested: n = 50000 bits, 38 errors.

BER\textsubscript{new} = 38/50000 = 7.6\ensuremath{\times}10\textsuperscript{-4}

H\textsubscript{0}: p = 0.001, H\textsubscript{1}: p \textless{} 0.001 (improvement), \ensuremath{\alpha} = 0.05

Z = (\ensuremath{\hat{p}} - p\textsubscript{0})/\ensuremath{\surd}(p\textsubscript{0}(1-p\textsubscript{0})/n) = (0.00076 - 0.001)/\ensuremath{\surd}(0.001\ensuremath{\times}0.999/50000) = -0.00024/0.000141 = -1.70

p-value = P(Z \textless{} -1.70) = 0.0446 \textless{} 0.05 \ensuremath{\to} \textbf{Reject H\textsubscript{0}.}

Conclusion: Evidence supports that the new algorithm has lower BER.

\end{PSbox}

\begin{PSbox}[unbreakable]{ex}{Example 4: Has Noise Power Increased?}

Historical: \ensuremath{\sigma}\textsubscript{0}\textsuperscript{2} = 0.04 V\textsuperscript{2}. New measurements: n = 25, S\textsuperscript{2} = 0.058 V\textsuperscript{2}.

H\textsubscript{0}: \ensuremath{\sigma}\textsuperscript{2} = 0.04, H\textsubscript{1}: \ensuremath{\sigma}\textsuperscript{2} \textgreater{} 0.04, \ensuremath{\alpha} = 0.05

\ensuremath{\chi}\textsuperscript{2} = (24)(0.058)/0.04 = 34.8

\ensuremath{\chi}\textsuperscript{2}\textsubscript{0.05, 24} = 36.415. Since 34.8 \textless{} 36.415 \ensuremath{\to} \textbf{Fail to reject H\textsubscript{0}.}

\end{PSbox}

\PSsection{15.12}{MATLAB Examples}

\begin{PSbox}[unbreakable]{ex}{Example 1: One-Sample t-test}

\tcbset{PScodemode/.style={unbreakable}}

\begin{Shaded}
\begin{Highlighting}[]
\CommentTok{\%\% Is the sensor biased? One{-}sample t{-}test}
\VariableTok{data} \OperatorTok{=}\NormalTok{ [}\FloatTok{0.012}\OperatorTok{,} \OperatorTok{{-}}\FloatTok{0.005}\OperatorTok{,} \FloatTok{0.031}\OperatorTok{,} \FloatTok{0.008}\OperatorTok{,} \OperatorTok{{-}}\FloatTok{0.012}\OperatorTok{,} \FloatTok{0.022}\OperatorTok{,} \FloatTok{0.015}\OperatorTok{,} \OperatorTok{...}
        \FloatTok{0.003}\OperatorTok{,} \FloatTok{0.028}\OperatorTok{,} \OperatorTok{{-}}\FloatTok{0.001}\OperatorTok{,} \FloatTok{0.018}\OperatorTok{,} \FloatTok{0.009}\OperatorTok{,} \FloatTok{0.025}\OperatorTok{,} \OperatorTok{{-}}\FloatTok{0.008}\OperatorTok{,} \OperatorTok{...}
        \FloatTok{0.014}\OperatorTok{,} \FloatTok{0.007}\OperatorTok{,} \FloatTok{0.020}\OperatorTok{,} \FloatTok{0.011}\OperatorTok{,} \OperatorTok{{-}}\FloatTok{0.003}\OperatorTok{,} \FloatTok{0.019}\NormalTok{]}\OperatorTok{;}

\VariableTok{mu\_0} \OperatorTok{=} \FloatTok{0}\OperatorTok{;}  \CommentTok{\% Hypothesized mean (no bias)}
\NormalTok{[}\VariableTok{h}\OperatorTok{,} \VariableTok{p\_value}\OperatorTok{,} \VariableTok{ci}\OperatorTok{,} \VariableTok{stats}\NormalTok{] }\OperatorTok{=} \VariableTok{ttest}\NormalTok{(}\VariableTok{data}\OperatorTok{,} \VariableTok{mu\_0}\NormalTok{)}\OperatorTok{;}

\VariableTok{fprintf}\NormalTok{(}\SpecialStringTok{\textquotesingle{}Test result: H = \%d (1=reject H₀)\textbackslash{}n\textquotesingle{}}\OperatorTok{,} \VariableTok{h}\NormalTok{)}\OperatorTok{;}
\VariableTok{fprintf}\NormalTok{(}\SpecialStringTok{\textquotesingle{}p{-}value = \%.4f\textbackslash{}n\textquotesingle{}}\OperatorTok{,} \VariableTok{p\_value}\NormalTok{)}\OperatorTok{;}
\VariableTok{fprintf}\NormalTok{(}\SpecialStringTok{\textquotesingle{}t{-}statistic = \%.3f, df = \%d\textbackslash{}n\textquotesingle{}}\OperatorTok{,} \VariableTok{stats}\NormalTok{.}\VariableTok{tstat}\OperatorTok{,} \VariableTok{stats}\NormalTok{.}\VariableTok{df}\NormalTok{)}\OperatorTok{;}
\VariableTok{fprintf}\NormalTok{(}\SpecialStringTok{\textquotesingle{}95\%\% CI for μ: [\%.4f, \%.4f]\textbackslash{}n\textquotesingle{}}\OperatorTok{,} \VariableTok{ci}\NormalTok{)}\OperatorTok{;}
\end{Highlighting}
\end{Shaded}

\end{PSbox}

\begin{PSbox}[unbreakable]{ex}{Example 2: Two-Sample t-test}

\tcbset{PScodemode/.style={unbreakable}}

\begin{Shaded}
\begin{Highlighting}[]
\CommentTok{\%\% Compare two algorithms: is there a significant difference?}
\VariableTok{BER\_algo1} \OperatorTok{=}\NormalTok{ [}\FloatTok{0.0012}\OperatorTok{,} \FloatTok{0.0008}\OperatorTok{,} \FloatTok{0.0015}\OperatorTok{,} \FloatTok{0.0010}\OperatorTok{,} \FloatTok{0.0013}\OperatorTok{,} \OperatorTok{...}
             \FloatTok{0.0011}\OperatorTok{,} \FloatTok{0.0009}\OperatorTok{,} \FloatTok{0.0014}\OperatorTok{,} \FloatTok{0.0012}\OperatorTok{,} \FloatTok{0.0010}\NormalTok{]}\OperatorTok{;}
\VariableTok{BER\_algo2} \OperatorTok{=}\NormalTok{ [}\FloatTok{0.0007}\OperatorTok{,} \FloatTok{0.0009}\OperatorTok{,} \FloatTok{0.0006}\OperatorTok{,} \FloatTok{0.0008}\OperatorTok{,} \FloatTok{0.0005}\OperatorTok{,} \OperatorTok{...}
             \FloatTok{0.0008}\OperatorTok{,} \FloatTok{0.0007}\OperatorTok{,} \FloatTok{0.0006}\OperatorTok{,} \FloatTok{0.0009}\OperatorTok{,} \FloatTok{0.0007}\NormalTok{]}\OperatorTok{;}

\NormalTok{[}\VariableTok{h}\OperatorTok{,} \VariableTok{p\_value}\OperatorTok{,} \VariableTok{ci}\OperatorTok{,} \VariableTok{stats}\NormalTok{] }\OperatorTok{=} \VariableTok{ttest2}\NormalTok{(}\VariableTok{BER\_algo1}\OperatorTok{,} \VariableTok{BER\_algo2}\NormalTok{)}\OperatorTok{;}
\VariableTok{fprintf}\NormalTok{(}\SpecialStringTok{\textquotesingle{}Two{-}sample t{-}test:\textbackslash{}n\textquotesingle{}}\NormalTok{)}\OperatorTok{;}
\VariableTok{fprintf}\NormalTok{(}\SpecialStringTok{\textquotesingle{}  H = \%d (1=reject H₀: means are equal)\textbackslash{}n\textquotesingle{}}\OperatorTok{,} \VariableTok{h}\NormalTok{)}\OperatorTok{;}
\VariableTok{fprintf}\NormalTok{(}\SpecialStringTok{\textquotesingle{}  p{-}value = \%.4e\textbackslash{}n\textquotesingle{}}\OperatorTok{,} \VariableTok{p\_value}\NormalTok{)}\OperatorTok{;}
\VariableTok{fprintf}\NormalTok{(}\SpecialStringTok{\textquotesingle{}  t{-}stat = \%.3f, df = \%.1f\textbackslash{}n\textquotesingle{}}\OperatorTok{,} \VariableTok{stats}\NormalTok{.}\VariableTok{tstat}\OperatorTok{,} \VariableTok{stats}\NormalTok{.}\VariableTok{df}\NormalTok{)}\OperatorTok{;}
\VariableTok{fprintf}\NormalTok{(}\SpecialStringTok{\textquotesingle{}  Mean diff: \%.4f\textbackslash{}n\textquotesingle{}}\OperatorTok{,} \VariableTok{mean}\NormalTok{(}\VariableTok{BER\_algo1}\NormalTok{) }\OperatorTok{{-}} \VariableTok{mean}\NormalTok{(}\VariableTok{BER\_algo2}\NormalTok{))}\OperatorTok{;}
\end{Highlighting}
\end{Shaded}

\end{PSbox}

\begin{PSbox}[unbreakable]{ex}{Example 3: Power Analysis}

\tcbset{PScodemode/.style={unbreakable}}

\begin{Shaded}
\begin{Highlighting}[]
\CommentTok{\%\% Power: How likely are we to detect a true difference?}
\VariableTok{mu\_0} \OperatorTok{=} \FloatTok{0}\OperatorTok{;}          \CommentTok{\% Null hypothesis mean}
\VariableTok{mu\_true} \OperatorTok{=} \FloatTok{0.015}\OperatorTok{;}   \CommentTok{\% True mean (sensor has bias of 15mV)}
\VariableTok{sigma} \OperatorTok{=} \FloatTok{0.03}\OperatorTok{;}      \CommentTok{\% Known std dev}
\VariableTok{alpha} \OperatorTok{=} \FloatTok{0.05}\OperatorTok{;}

\VariableTok{n\_values} \OperatorTok{=} \FloatTok{5}\OperatorTok{:}\FloatTok{5}\OperatorTok{:}\FloatTok{100}\OperatorTok{;}
\VariableTok{power} \OperatorTok{=} \VariableTok{zeros}\NormalTok{(}\VariableTok{size}\NormalTok{(}\VariableTok{n\_values}\NormalTok{))}\OperatorTok{;}

\KeywordTok{for} \VariableTok{i} \OperatorTok{=} \FloatTok{1}\OperatorTok{:}\VariableTok{length}\NormalTok{(}\VariableTok{n\_values}\NormalTok{)}
    \VariableTok{n} \OperatorTok{=} \VariableTok{n\_values}\NormalTok{(}\VariableTok{i}\NormalTok{)}\OperatorTok{;}
    \CommentTok{\% Critical value for two{-}sided test}
    \VariableTok{z\_crit} \OperatorTok{=} \VariableTok{norminv}\NormalTok{(}\FloatTok{1}\OperatorTok{{-}}\VariableTok{alpha}\OperatorTok{/}\FloatTok{2}\NormalTok{)}\OperatorTok{;}
    \CommentTok{\% Power = P(reject | mu\_true)}
    \CommentTok{\% Reject when |Z| \textgreater{} z\_crit, where Z \textasciitilde{} N((mu\_true{-}mu\_0)/(sigma/sqrt(n)), 1)}
    \VariableTok{ncp} \OperatorTok{=}\NormalTok{ (}\VariableTok{mu\_true} \OperatorTok{{-}} \VariableTok{mu\_0}\NormalTok{) }\OperatorTok{/}\NormalTok{ (}\VariableTok{sigma}\OperatorTok{/}\VariableTok{sqrt}\NormalTok{(}\VariableTok{n}\NormalTok{))}\OperatorTok{;}  \CommentTok{\% Non{-}centrality parameter}
    \VariableTok{power}\NormalTok{(}\VariableTok{i}\NormalTok{) }\OperatorTok{=} \FloatTok{1} \OperatorTok{{-}} \VariableTok{normcdf}\NormalTok{(}\VariableTok{z\_crit} \OperatorTok{{-}} \VariableTok{ncp}\NormalTok{) }\OperatorTok{+} \VariableTok{normcdf}\NormalTok{(}\OperatorTok{{-}}\VariableTok{z\_crit} \OperatorTok{{-}} \VariableTok{ncp}\NormalTok{)}\OperatorTok{;}
\KeywordTok{end}

\VariableTok{figure}\OperatorTok{;}
\VariableTok{plot}\NormalTok{(}\VariableTok{n\_values}\OperatorTok{,} \VariableTok{power}\OperatorTok{,} \SpecialStringTok{\textquotesingle{}b{-}o\textquotesingle{}}\OperatorTok{,} \SpecialStringTok{\textquotesingle{}LineWidth\textquotesingle{}}\OperatorTok{,} \FloatTok{2}\NormalTok{)}\OperatorTok{;}
\VariableTok{xlabel}\NormalTok{(}\SpecialStringTok{\textquotesingle{}Sample Size n\textquotesingle{}}\NormalTok{)}\OperatorTok{;} \VariableTok{ylabel}\NormalTok{(}\SpecialStringTok{\textquotesingle{}Power (1 {-} β)\textquotesingle{}}\NormalTok{)}\OperatorTok{;}
\VariableTok{title}\NormalTok{(}\VariableTok{sprintf}\NormalTok{(}\SpecialStringTok{\textquotesingle{}Power vs Sample Size (true bias = \%.0f mV)\textquotesingle{}}\OperatorTok{,} \VariableTok{mu\_true}\OperatorTok{*}\FloatTok{1000}\NormalTok{))}\OperatorTok{;}
\VariableTok{grid} \VariableTok{on}\OperatorTok{;}
\VariableTok{yline}\NormalTok{(}\FloatTok{0.8}\OperatorTok{,} \SpecialStringTok{\textquotesingle{}r{-}{-}\textquotesingle{}}\OperatorTok{,} \SpecialStringTok{\textquotesingle{}Desired Power = 0.8\textquotesingle{}}\NormalTok{)}\OperatorTok{;}
\VariableTok{fprintf}\NormalTok{(}\SpecialStringTok{\textquotesingle{}Need n ≈ \%d for 80\%\% power\textbackslash{}n\textquotesingle{}}\OperatorTok{,} \VariableTok{n\_values}\NormalTok{(}\VariableTok{find}\NormalTok{(}\VariableTok{power} \OperatorTok{\textgreater{}=} \FloatTok{0.8}\OperatorTok{,} \FloatTok{1}\NormalTok{)))}\OperatorTok{;}
\end{Highlighting}
\end{Shaded}

\end{PSbox}

\PSsection{15.13}{Practice Problems}

\begin{PSbox}[unbreakable]{prob}{Problem 1}

A manufacturing spec requires mean resistance = 100\ensuremath{\Omega}. A sample of 16 resistors gives \ensuremath{\bar{X}} = 101.5\ensuremath{\Omega}, S = 3\ensuremath{\Omega}. At \ensuremath{\alpha} = 0.05, is there evidence the process has drifted?

\PSsolution{} H\textsubscript{0}: \ensuremath{\mu}=100, H\textsubscript{1}: \ensuremath{\mu}\ensuremath{\neq}100. T = (101.5-100)/(3/4) = 2.0. t\textsubscript{0.025,15} = 2.131. \textbar{}T\textbar{}=2.0 \textless{} 2.131 \ensuremath{\to} Fail to reject. p-value \ensuremath{\approx} 0.064 \textgreater{} 0.05.

\end{PSbox}

\begin{PSbox}[unbreakable]{prob}{Problem 2}

Average packet delay was 5ms. After network upgrade, 40 measurements give \ensuremath{\bar{X}} = 4.6ms, S = 1.2ms. Test if delay decreased at \ensuremath{\alpha} = 0.01.

\PSsolution{} H\textsubscript{0}: \ensuremath{\mu}=5, H\textsubscript{1}: \ensuremath{\mu}\textless{}5. T = (4.6-5)/(1.2/\ensuremath{\surd}40) = -0.4/0.190 = -2.11. t\textsubscript{0.01,39} \ensuremath{\approx} -2.426. T = -2.11 \textgreater{} -2.426 \ensuremath{\to} Fail to reject at 1\% level. (Would reject at 5\% since t\textsubscript{0.05,39}\ensuremath{\approx}-1.685.)

\end{PSbox}

\begin{PSbox}[unbreakable]{prob}{Problem 3}

Explain why ``p = 0.04 means there is only a 4\% chance H\textsubscript{0} is true'' is WRONG.

\PSsolution{} The p-value is P(data this extreme \textbar{} H\textsubscript{0} true), NOT P(H\textsubscript{0} true \textbar{} data). It conditions on H\textsubscript{0} being true and asks about data extremity. To get P(H\textsubscript{0} true \textbar{} data) would require Bayesian analysis with a prior. The p-value is a property of the data-generating process under H\textsubscript{0}, not a posterior probability about H\textsubscript{0}.

\emph{Next Module: {Module 16 — Engineering Statistical Applications}}

\end{PSbox}
